\subsection{Classical early-warning indicators}
We compared trajectory-based warnings with four classical statistical indicators: lag-one autocorrelation, variance, an AR(1)-based relaxation-rate proxy, and a low-to-high-frequency power ratio \cite{Dakos2012EWS}. The input was the causal robust amplitude envelope, normalized by its positive median in the fitting prefix and sampled every 0.05 s. Each indicator was calculated from a trailing $5T$ window after local linear detrending. Windows of $3T$ and $7T$, and a $5T$ window applied to $R_A$, were sensitivity analyses.
For detrended samples $v_1,\ldots,v_N$, we used $AC_1=\sum_{k=2}^N v_kv_{k-1}/\sum_{k=1}^N v_k^2$ and sample variance with denominator $N-1$. The fitted coefficient was $\phi=\sum_{k=2}^N v_kv_{k-1}/\sum_{k=1}^{N-1}v_k^2$. We defined $\kappa=-\log(\phi)/\Delta t$ only for $0<\phi<1$, with $\Delta t=0.05$ s; other windows provided no recovery-rate warning evidence. This is an AR(1) proxy, not a direct measurement of the physiological recovery rate. Spectral reddening was the ratio of summed Hann-periodogram power in $0<f\leq0.5/T$ to that in $0.5/T<f\leq2/T$.
Warning directions were fixed: increasing autocorrelation, variance and spectral ratio, and decreasing $\kappa$. Each recording's threshold candidates were fitting-prefix quantiles. A single quantile level per indicator was selected across records on validation data by the harmonic mean of event coverage and complete-follow-up alarm precision. This selection score is not conventional event F1 because repeated warnings can associate with one event. Selected thresholds were frozen before test scoring. All methods shared three-anchor entry confirmation, 0.10-s exit confirmation, the evaluation time grid, and the $2T$ actual-issuance association window. Scalar indicators were not required to supply fictitious future trajectories. The existing operator and persistence working points were retained.
We aligned causal indicator histories to the 88 test reference onsets and estimated event-wise Kendall trends and late-minus-early changes. Pointwise intervals and paired method differences used 2,000 resamples of the seven participant clusters. Intervals condition on selected thresholds. Threshold sweeps on the test set describe the coverage--burden trade-off and were not used to choose a deployment setting. Separate 20\%, 30\% and 40\% active-time targets were selected on validation data and evaluated at their actual test burdens.
